\section{Finite cochain models and determinant comparisons}\label{sec:foundations}

The determinant arguments use cohomology at moving finite levels, together
with local conditions, cup products, and distinguished classes. We first
construct finite free models that retain these maps under passage to a
limit. We then prove the evaluation and determinant comparisons used to
change local conditions. The final two subsections compute the central
Kummer volumes and give the bounded-series tests needed for specialization.

We use local Tate duality and global Poitou--Tate duality, including compact
support with Tate modification at real places; see \cite{MilneADT}. The
mapping-fiber and Bockstein language is that of Selmer
complexes~\cite{Nekovar06}. At two, the standing hypotheses in that reference
impose a totally imaginary base field; our ordinary-positive and Tate-real
comparisons over \(\mathbb Q\) are given explicitly below. The other standard
inputs are étale localization and purity, Shapiro's Lemma, and arithmetic
cohomological dimension and finiteness theorems.

A global support \(S\) for a number field \(F\) includes the primes over 2
and the archimedean places; a list just of finite allowed primes leaves the
archimedean places understood. Write \(G_{F,S}\) for the
unramified-outside-\(S\) quotient. Frobenius in a coefficient action means
arithmetic Frobenius. Translating that action to a weight on a CM translate,
group polynomial, or Euler product can therefore introduce an inverse
variable; the weights in each comparison use the same convention.

For finite or compact 2-primary coefficients \(M\), \(G(M)\) and \(L_v(M)\) denote the unrestricted global and local continuous cochain complexes for the indicated support. We also write \(C_{\rm glob}(M)\) and \(C_v(M)\). A Selmer complex with local maps \(U_v\to L_v\) is
\[
 C(U,M)=\mathrm{fib}\big(G(M)\longrightarrow\bigoplus_v Q_v\big),\qquad
                Q_v=\mathrm{cone}(U_v\to L_v);
\]
we use cohomological mapping fibers (cone shifted by \(-1\)). “Strict” and “full” usually mean \(U_v=0\) and \(L_v\), respectively, at the indicated place only. “Finite” will refer to a Kummer or an unramified/Frobenius condition as specified. We omit terms imposing no condition. A separate positive complex over \(\mathbb Q\) uses the mapping fiber of restriction to \textbf{ordinary} \(C^\bullet(\mathbb R,M)\). Complex places do not contribute local conditions. By dual on coefficients we mean Tate dual (linear dual with twist (1), or compact/discrete duals as appropriate). Write \(P^\vee\) on complexes for the derived linear dual with \emph{no} shift unless specified.

\subsection{Finite models with marked arithmetic maps}

Some parameters come from a fixed tower, while others come from cyclic
extensions whose conductors move. Both are treated at finite Artin
precision. For the moving parameters, the limit is taken on the matrices
of the arithmetic diagrams; it is not cohomology of a new infinite-conductor
Galois group.

For example, a cyclic generator of order \(2^m\), written \(1+u\), gives
the relation \((1+u)^{2^m}-1\). At the fixed precision
\((2^n,u^n)\) this relation vanishes for all sufficiently large \(m\).
Thus the growing cyclic group rings have the same finite precisions as
\(\mathbb Z_2[[u]]\). The following lemma carries the cochain maps, as
well as the coefficient rings, through this passage.

\begin{lemma}[Finite models and matrix limits]\label{lem:foundations-finite} Consider diagrams of the global, local, and mapping-fiber cochain complexes above over finite local Artin coefficient rings of 2-power characteristic, with modules finite free (and compatible lifts), in systems with fixed finite residue field and only boundedly many moving places. Assume the dimensions of the residual cohomology of each complex to be modeled are uniformly bounded per degree; if a bounded perfect model is used integrally, assume these dimensions vanish outside uniformly bounded degrees as well. Then the following constructions can be made by cochain contractions.

\begin{enumerate}
\item There are degreewise finite free minimal models with ordinary base change, bounded where indicated. Continuous actions over complete rings at a stage can be treated with compatible contractions. Tensoring a model of a free coefficient also computes cochains with its tensor by any fixed finite scalar module with trivial Galois action (at compact precision, or with completed finite module).
\item Suppose on a fixed nonprincipal ultrafilter on a sequence of stages the scalar precisions identify, by reduction, with every fixed finite Artin precision of a complete local limiting ring \(R\). Here and below arbitrary such increasing precisions cofinal with powers of the maximal ideal can be used. Transfer cochain diagrams at the stages and take matrix limits at all fixed precisions. The resulting finite free models retain maps, homotopies, exact triangles (with their identifications), specified cycles and cup-product/duality diagrams when marked as below. Any fixed continuous integral parameter substitution commutes with this construction. Fixed-group diagrams whose coefficient actions already stabilize mod each precision in this identification compute the actual cohomology and maps. One can extend a previously taken diagram by further moving-prime operations on the \emph{same} stages and ultrafilter, retaining previous models and maps; one can work to a slower cofinal precision for the extensions.
\end{enumerate}

\end{lemma}

\paragraph{Arithmetic hypotheses for the models.}
The concrete limit rings for arithmetic diagrams will be power series over finite \(2\)-adic integer rings, or over a fixed finite complete local algebra, possibly after base change. The number of moving primes counts distinct places, not ramification degrees.

On global complexes the residual bounds required here, for a fixed residual representation up to a bounded filtration by fixed ones, follow by a fixed finite extension killing those representations and Kummer theory there (and Hochschild--Serre or inflation bounds and duality). For instance valuations allowed in squareclasses increase by only a bounded number of coordinates; degree two is bounded by Poitou--Tate. The local bounds use the usual local theory; degrees above two globally contribute only real cohomology and restriction to all real places there is an isomorphism. In particular positive complexes with no other conditions have the requisite bounded range, and over a totally imaginary field the ordinary global complexes do too.

\begin{proof}
\emph{Contraction at finite precision.}
Continuous cochains with free coefficients at finite precision are flat-term (filtered colimit over finite clopen partitions); reduction or tensor of the terms is ordinary cochains. Over a local Artin ring flat modules are free: lift a residue basis, use nilpotence for surjectivity and then flatness for the kernel modulo the maximal ideal followed by nilpotence again. Residually split the vector-space complex into its cohomology with zero differential and a sum of disks (a differential isomorphism with bases in adjacent degrees), and lift the graded bases. Let \(d_0\) be the pure disk differential in them, \(i,p,h\) the inclusion, projection and contraction with \(h^2=hi=ph=0,\ 1-ip=d_0h+hd_0\). For the actual \(d=d_0+\epsilon\), use
\[
 \begin{aligned}
 I&=(1+h\epsilon)^{-1}i,& P&=p(1+\epsilon h)^{-1},\\
 H&=h(1+\epsilon h)^{-1},& d_{\rm small}&=p\epsilon(1+h\epsilon)^{-1}i.
 \end{aligned}
\]
These are the disk contraction formulas (multiplication using \(d_0\epsilon+\epsilon d_0+\epsilon^2=0\) gives \(dI=I d_{\rm small}, P d=d_{\rm small}P, PI=1, dH+Hd=1-IP\)); the series are finite in the nilpotent ideal. 

\emph{Compatible diagrams.}
Bases of continuous cochains at successive quotients can be lifted compatibly by surjectivity on functions; the resulting maps of inverse limits give the same identities. Transfers of maps insert \(I,P\); compositions have the homotopies from \(H\). Analogously transfer multilinear cup maps in the required degrees by tensoring the contractions. For a mapping cone/fiber one can use the cone/fiber of a transferred map with the induced homotopy equivalences, or its minimal contraction. Acyclic reductions of cones can be retained with contracting homotopies.

\emph{Matrix limits and fixed-group comparison.}
Ranks and all finite scalar entries used in these identities have simultaneous ultrafilter values, degree by degree and compatibly by reduction; zero regions of minimal models remain zero. For fixed-group comparison at each compact scalar quotient one can compare to a further fixed contraction for the ordinary cochains at that quotient by composing inclusions and projections. On cohomology the inverse isomorphisms are between finite groups and compatible across quotients by the actual cochain reductions. The same comparisons commute with the marked maps (or one transfers the comparison homotopies on the fixed models too). Passing to compact inverse limits has no \(\lim^1\) issue, by finiteness/compactness. 

Extensions of a diagram use the old contractions on the needed old objects on ultrafilter-large sets at each precision (matrix bases can be padded or stabilized per degree). One may choose a slower single cofinal working precision on these sets. At no point does one define a new infinite-conductor profinite Galois group and posit Poitou--Tate on it.

\end{proof}

\subsection{Localization and compatible duality}

\paragraph{Change of support.}
Over a totally imaginary field the arithmetic \(S\)-integer scheme computes \(G_{F,S}\)-cohomology for our lisse systems. Indeed the 2-primary comparison is checked residually after finite covers trivializing coefficients; in the universal such cover higher cohomology of \(\mathbb F_2=\mu_2\) vanishes. Degree one is the torsor criterion, and in degree two use Picard modulo 2 and Brauer. Brauer classes die by even local degrees eventually at all of \(S_f\) in the cyclotomic tower (also at odd residue cardinalities, which have infinite order), then Picard classes by Hilbert principalization. Higher degrees vanish already over imaginary layers (use localization to the generic point and purity, cohomological dimension, and in degree three surjectivity of Brauer residues away from \(S\) by Brauer reciprocity since \(S_f\ne\varnothing\)). 

Thus allowing additional odd places where a coefficient is unramified changes the full complex by the sum of local singular quotients, i.e. local cochains modulo the inflated
\[
                         [M\xrightarrow{\mathrm{Fr}-1}M]\qquad(0,1).
\]
Inflation into the new global problem imposing just these unramified conditions there is the quasi-isomorphism (étale localization and local excision); fixed conditions elsewhere commute with it. The same works at compact precisions and limits. At unramified coefficients a singular block can be computed as the residue cochains on \(M(-1)\) shifted by \(-1\). More generally local computations at odd places discard the exact-invariants pro-(prime-to-2) kernel of inertia, then use procyclic \(2\)-inertia and residue cohomology. The inertia differential in degrees 0,1 uses generator minus one; under a power \(z\mapsto z^b\) the change on the generator cochain uses \((z^b-1)/(z-1)\) in completed procyclic notation. These are finite-model formulas (the usual two-term procyclic resolution, followed by the residue Frobenius fiber), so when the conjugating power tends to 1 the change term tends to 1 on finite precisions with pro-2 action. Unramified conditions after inverting 2 at odd places with fixed possibly ramified Tate action will be described separately.

\paragraph{Cup products and traces.}
On an imaginary base retain full global, full finite local \(L=\bigoplus L_v\) and compact \(G_c=\mathrm{fib}(G\to L)\) diagrams. Cup/invariant give a perfect pairing of \(G_c\) with the dual coefficient's unrestricted \(G'\) of shift \(-3\), and of \(L,L'\) of shift \(-2\); the connecting \(L\to G_c[1]\) is adjoint to localization, up to orientation of fibers. These are standard Poitou--Tate/local duality over the residue field, and hold over the scalar rings with free coefficient models as well. 

More explicitly transfer the cup maps, invariant and compact trace (to scalars in degrees two locally and three globally), and the reciprocity/adjunction diagram including its homotopies. The traces on scalar cochains with twist (1) can be taken at finite precision first: top cohomology is the scalar ring by the invariant or compact invariant theorem, compatibly and linearly under reduction (dually the constant degree-zero invariants of the finite dual); higher cohomology vanishes. Thus on minimal models the top degree term for each trace is one copy with zero incoming differential and the trace gives a morphism there. 

The compact map and local invariants agree on boundaries in this diagram by the ordinary compact-support triangle. The resulting cup morphisms lift residual perfectness to equivalences by the finite models. Equivalently, for the full compact/global adjunction use the cone cochain cup into compact scalar cochains, and trace; restriction on the unrestricted factor then pairs with local boundary by the local trace. These equivalences and the triangle adjunction survive the matrix limit and base changes. One may incorporate conjugate field transport in the identification of the second coefficient with a Tate dual (e.g. an anticyclotomic scalar twist); in that convention exact complementary conditions must be tested at the corresponding conjugate places.

\begin{lemma}[Complementary local conditions]\label{lem:foundations-duality}
If paired perfect local maps \(U\to L,\ U'\to L'\) (possibly only over a localized test ring after the limit) have a paired restriction nullhomotopy and induce \(L'/U'\simeq U^\vee[-2]\), the two Selmer complexes are perfect duals with shift \(-3\). \end{lemma}

\begin{proof}
Dualize \(G_c\to C(U)\to U\) and use the localization adjunction to identify the dual fiber. The comparison boundaries for nested local conditions with compatible orthogonality data are localization adjoints by the same triangles. The local equivalence can be checked on residue fields by perfectness; when the local pieces are split free in cohomology, the invariant restriction morphism is tested just in total degree two. We specify nullhomotopies or their existence from degree/amplitude or unramified factorizations when used; exact Selmer orthogonality on a DVR does not follow just from two generic dimensions.

\end{proof}

\paragraph{The real place.}
Over \(\mathbb Q\), integral calculations without a perfectness assertion for full dual pairs use ordinary positive models and finite Poitou--Tate (modified Tate at the real place in that duality). In 2-inverted comparisons the real term of the positive-complex triangle contributes just the invariant module, in degree one in the determinant of the fiber. Ordinary global complexes then have the imaginary-base comparisons above: one can first include diagrams by restriction/induction to \(\mathbb Q(i)\), whose ramification is allowed, and project by the two maps after inverting 2. Restriction/corestriction and Shapiro are cochain identities up to homotopy, transferred in the finite diagrams, so the ordinary models, even if originally unbounded at real cohomology, become perfect direct summands there. Alternatively use the splitting of induced coefficient morphisms with denominator 2. This proves the duality and unramified inflation comparisons in that setting; the averaging scalar in adjunction is invertible there. 

A different integral use over \(\mathbb Q\) is with Shapiro coefficients from a quadratic imaginary field, when ramification of that field is allowed. The modified real term of such a coefficient is contractible even integrally. The paired Shapiro cup pushed by trace to scalar cochains uses local/compact invariant \textbf{compatible with corestriction}, not twice the invariants over the imaginary base. Compact scalar cochains here use modified Tate localization at infinity. They can still be used degreewise for their degree-three trace as above (top cohomology by finite Poitou--Tate); the bounded dual pair and comparison maps for the induced coefficients equivalently identify by Shapiro. At discriminant primes newly allowed one retains the unramified Kummer/cochain conditions over the imaginary field as appropriate.

\paragraph{Character pushforward.}
Twisted classes obtained by a coefficient character map on an induced group-ring coefficient are unnormalized traces. Concretely, restricting the resulting corestriction up to a field trivializing that character gives the sum over translates with the character weights, not division by the group order. This applies at finite precisions, including when 2 is not inverted. Conversely at a fixed characteristic-zero character one may compare local conditions using its projected summand before further inverse limits if projector denominators are then bounded.

\subsection{Evaluations and Frobenius tests}

\begin{lemma}[Detection by evaluations]\label{lem:foundations-evaluation}
For the limiting unrestricted global \(H^1\), after scalar specialization even to height-one residue fields, there is an injection by evaluations into abstract crossed homomorphism classes on \(\mathcal G_F=\prod_i G_F\), where \(i\) runs over the stages. The action on a fixed-rank module is given by entrywise ultrafilter limits before specialization. \end{lemma}

\begin{proof}
Retain the model inclusions into the actual global cochains. Evaluations
on every sequence of group elements, pairs of sequences, and higher tuples
have simultaneous matrix limits. They preserve the cochain equations, so
degree-one cycles give crossed homomorphisms on \(\mathcal G_F\).

It remains to show that evaluation detects their cohomology classes after
specialization. At each residual stage a bounded finite list of group
elements detects invariants and whether a model cocycle is a single
coboundary. These are linear equations on bounded-dimensional spaces.
The resulting evaluation map to
\[
                         E=[M\longrightarrow M^{\rm list}]
                         \qquad\text{in degrees }0,1
\]
is therefore an isomorphism on residual \(H^0\) and injective on residual
\(H^1\). For its unshifted mapping cone \(K\), writing \(k\) for the
residue field, the long exact sequence gives
\[
                         H^{-1}(K\otimes k)=H^0(K\otimes k)=0.
\]
The complexes start in degree zero, so a degreewise minimal model of
\(K\) has no terms in degrees at most zero. This remains true after every
field base change. The same long exact sequence now proves the required
\(H^1\)-injection at the specialized field. No upper boundedness of the
ordinary real cochains is needed for this low-degree argument.

Conjugation and comparison identities on evaluations hold on classes by
the transferred homotopies. A Selmer \(H^1\) likewise injects whenever its
local maps include full \(H^0\) and inject in \(H^1\) at the field being
tested.
\end{proof}

\paragraph{Detection on a kernel.}
Let \(\mathcal H\) be a normal joint kernel acting trivially on the
coefficients. Restriction to \(\mathcal H\) gives additive, equivariant
evaluations. Inflation--restriction shows that these detect classes if
\(H^1(\mathcal G_F/\mathcal H,M)=0\). One sufficient condition is a central
scalar whose difference from one is invertible, by the crossed-homomorphism
identity.

\paragraph{The open-image criterion.}
A second criterion applies in characteristic zero. Suppose the quotient
contains, as a normal subgroup, the product
\[
                   \mathcal S=\prod_i\mathrm{SL}_2^b(\mathbb Z_2)
\]
of a fixed deep principal congruence subgroup, acting on the standard
plane by the limiting matrices, with all scalar twists trivial on
\(\mathcal S\). Then
\[
                         H^0(\mathcal S,M)=H^1(\mathcal S,M)=0,
\]
where \(H^1\) denotes even abstract group cohomology. These vanishings imply
the required vanishing on the quotient.

To prove them, choose the constant diagonal
\(D=\operatorname{diag}(m,m^{-1})\), with an integer \(m>1\) sufficiently
close to one. The matrix \(D-1\) is invertible. Adjust a cocycle to vanish
at \(D\); commutation with \(D\) then makes it vanish on every sequence
of diagonal elements. For an upper-unipotent sequence \(U\), conjugation
by \(D\) gives the integer power \(U^{m^2}\). The crossed identity gives
\[
             \left(D-\sum_{j=0}^{m^2-1}U^j\right)x(U)=0.
\]
This triangular matrix is invertible: its diagonal entries are
\(m-m^2\) and \(m^{-1}-m^2\). Using \(D^{-1}\) gives the same conclusion
for lower unipotents. Gaussian factorization into lower-unipotent,
diagonal, and upper-unipotent factors works termwise inside the same
principal congruence subgroup. Hence it works for arbitrary sequences,
and the cocycle vanishes on \(\mathcal S\).

This subgroup is available for an open elliptic Tate image when the
competing data have uniformly bounded derived length over a fixed finite
field extension. A sufficiently iterated closed derived subgroup kills
those data and contains a deep \(\mathrm{SL}_2\) in the Tate image, by
Lie brackets first in the open matrix group and then in
\(\mathfrak{sl}_2\). The depth can be chosen uniformly at the stages.

\paragraph{Joint evaluations and Frobenius selection.}
On a simple coefficient space with scalar endomorphisms, the joint evaluations of linearly independent classes detected by that restriction span a full tuple space. A proper span, by equivariance and semisimplicity of a sum of copies of the simple, would impose a scalar linear relation. In characteristic zero additive evaluations then permit simultaneous nonidentically-zero polynomial tests (integer combinations). Even in characteristic two, if full span gives rank two somewhere for a \emph{linear} pair-coordinate test, additivity gives rank two on the kernel evaluations: a homogeneous quadratic vanishing on an additive subgroup vanishes on its field span by polarization. 

Finally by Chebotarev a Frobenius sequence can agree with a selected sequence of elements on all old evaluation matrices and other designated finite continuous data to increasing precision: continuous functions at each Artin precision factor through finite sets on finite quotients, and one uses a chosen representative over the prime. In inert-prime tests of squares one uses the desired coset and takes normal closures of the bounded-precision data over the base. Additional splitting and determinant prescriptions will need actual compatibility, which will be checked in each use. This procedure preserves old evaluations when enlarging support (with unramified initial condition at a new place via the localization comparisons).

\subsection{Determinant volumes and complementary changes}

We use the determinant-line formalism for perfect complexes
\cite{KnudsenMumford76}, with the inverse convention specified below.
The complementary-condition switches are related to finite--singular
comparisons in \cite[Section~1.2]{MazurRubin04}; the following proofs
retain their exact local determinant volumes.

Use the line
\[
 \mathcal D(C)=\bigotimes_i(\det C^i)^{(-1)^{i+1}}
\]
for a perfect complex (the inverse of the cohomological determinant), and its canonical exact-triangle and cohomology identifications, ignoring harmless determinant signs. For a DVR with valuation \(v\), \(d(C,x)\) is the valuation of the coefficient of a generic tensor \(x\) relative to a generator of \(\mathcal D(C)\); an infinite value is allowed for zero. A torsion elementary divisor in cohomological degree \(i\) of length \(b\) contributes \((-1)^{i+1}b\), in uniformizer units of value one, to \(d\) on free cohomology bases. For a square \([F\xrightarrow{\delta}G]\) in degrees 1,2 with finite-length cohomology, \(d(C,1)=-v(\det\delta)\). These statements follow immediately by splitting into elementary-divisor blocks. For a self-duality of shift \(-3\) pairing degrees 1 and 2 of rank one generically, \((Y,Y^\vee)\) will denote the determinant tensor of a degree-one class \(Y\) with its pairing functional on degree two.

\begin{lemma}[The strict--finite determinant comparison]\label{lem:foundations-D}
Comparisons of local conditions \(U\subset V\) give a triangle \(C(U)\to C(V)\to V/U\). Suppose at a split pair of places one complex \(C_*\) uses strict at \(w\), full at \(\bar w\); \(C_F\) uses exact orthogonal \(U_w,U_{\bar w}\), keeping other exact orthogonal conditions, with conjugate duality (so \(Q_{\bar w}\simeq U_w^\vee[-2]\)). Assume perfect comparisons including duals over our DVR, and generic \(U_w\) a line in degree one only. Write \(C_0\) for strict at \(w\), \(U_{\bar w}\) at \(\bar w\). We have
\[
 C_0\to C_F\to U_w,\qquad C_0\to C_*\to Q_{\bar w}.
\]
Over the fraction field suppose \(C_*\) acyclic and a class \(Y\in H^1(C_F)\) has nonzero localization in \(H^1(U_w)\). The triangles show \(C_F\) has only a line in each of degrees 1,2 and degree one maps isomorphically onto that localization line. Conversely these line assertions force strict acyclicity (using duality: the boundary from \(Q_{\bar w}\) is then nonzero by adjunction). Write \(\operatorname{loc}_wY=p e\) for a fraction-field basis \(e\). Exactly
\[
 d(C_F,(Y,Y^\vee))=d(C_*,1)+2\big(v(p)+d(U_w,e)\big).       \tag{D}\label{eq:D}
\]
\end{lemma}

\begin{proof}
For \(x=\partial(e^*)\in H^2(C_0)\) by the second triangle, \(x\) pairs after projection with \(Y\) by \(p\), by the same adjunction. Thus up to signs the two tensors compared via \(C_0\) give \(p^2 x^*\otimes e,\ x^*\otimes e^*\), respectively; the two local lattice volumes are dual. The computation works also for rationally imposed local conditions in compatible triangle volumes defined by assigning a DVR determinant lattice on \(U_w\) and the dual one on \(Q_{\bar w}\).

\end{proof}

\subsection{Line and square switches}

A switch changes the condition at one newly allowed place. Fix nested
local conditions \(U_-\subset U_+\), and intermediate conditions
\(U_f,U_s\) whose quotients over \(U_-\) give a direct-sum decomposition
of \(U_+/U_-\). Over the test DVR, assume these two quotients are free
and concentrated in degree one. Their coordinates are called finite
and transverse (pure singular), and denoted \(f,s\). Let
\(C_-,C_+,C_f,C_s\) be the four resulting global complexes.
The local conditions form the diagram
\[
\begin{tikzpicture}[baseline=(lower.base),>=Stealth,node distance=8mm and 16mm]
\node (lower) {\(U_-\)};
\node (finite) [above right=of lower] {\(U_f\)};
\node (singular) [below right=of lower] {\(U_s\)};
\node (upper) [below right=of finite] {\(U_+\)};
\draw[->] (lower) -- (finite);
\draw[->] (lower) -- (singular);
\draw[->] (finite) -- (upper);
\draw[->] (singular) -- (upper);
\end{tikzpicture}
\qquad
 U_+/U_-\simeq(U_f/U_-)\oplus(U_s/U_-).
\]

We assume complementary dualities: the upper problem is dual to the
lower problem for the dual coefficient, and conversely. The local
pairing is perfect between the two coordinate spaces, and each pure
condition is self-annihilating against its corresponding condition on
the dual side. The global complexes have amplitude \(1,2\) over the
generic field and each indicated fiber field. The local splitting and
complementary duality hypotheses hold at both fields; the lemmas distinguish
the further generic and fiber dimension assumptions.

In the applications, the lower condition contains all local \(H^0\),
and the upper condition contains both degree-one coordinates but no
local \(H^2\). Thus the lower condition is generally not the zero
complex. Keeping this degree-zero term is what makes the changed global
\(H^1\) inject into unrestricted cohomology.

\begin{lemma}[Line switch]\label{lem:foundations-line}
Assume that the local quotient ranks are \(1,1\). Suppose \(C_f\) has just one generic cohomology line in degree one. At the fiber, if \(h=\dim H^1(C_f)>1\), arrange nonzero finite evaluations for both a primitive fiber reduction of that generic line and some degree-one class for the old dual condition. Then the switch lowers the fiber dimension from \(h\) to \(h-1\), and \(C_s\) again has a single generic cohomology line in degree one. If \(Y,Y'\) are classes for the finite and transverse conditions and \(s(Y')=f(Y)\) under an isomorphism of the two local lines of unit determinant volume, then
\[
 d(C_f,Y)=d(C_s,Y').
\]
If successive switches carry integral classes to a fiber of dimension one, this common coordinate has nonnegative valuation.
\end{lemma}

\begin{proof}
Then the generic finite evaluation is nonzero and \(C_-\) generically acyclic. On the fiber no transverse contribution can enter (it must annihilate the nonzero finite opposite evaluation by reciprocity), so \(H^1(C_s)\) has dimension \(h-1\) there by the triangles. Generically the triangle from lower gives again a line in \(C_s\), evaluating isomorphically on the local singular line.

If classes \(Y,Y'\) in \(C_f,C_s\) have \(s(Y')\) equal to \(f(Y)\) via a unit-volume line isomorphism, then
\[
                          d(C_f,Y)=d(C_s,Y')
\]
by the two triangles from \(C_-\) (same tensor 1 there). Extra unchanged terms with a fixed trivialization need not be included in this notation.

Thus in a sequence of such switches with integral classes at the test DVR and nonzero generic evaluations, when \(h=1\) the coordinate is integral (the complex is then just a free line by the minimal fiber model), and gives nonnegativity also backwards.
\end{proof}

\begin{lemma}[Square switch]\label{lem:foundations-square}
Assume that the local quotient ranks are \(2,2\), in the self-dual global setting with exact lower/upper duality. The localization image for upper in the four coordinates is self-annihilating: its annihilator is the kernel of the boundary into \(H^2(C_-)\), by adjunction, hence is the same image. Suppose \(C_f\) has generic rank one in degree one with a class \(Y\) and nonzero \(f(Y)=v_0\); then \(H^1(C_-)\) there vanishes. Suppose \(Y'\) for \(C_s\) has \(s(Y')=J v_0\) up to a unit, where \(J\) is an integral involution on a common coordinate plane, and the cross pairing of either plane with the other is \(\langle x,Jy\rangle\) up to a unit using a unimodular alternating plane form (interchanging the variables or transporting by \(J\) gives the same condition up to units). Then \(Y,Y'\) are a generic basis in upper by the half-dimension assertion, \(C_s\) has again generic rank one in degrees 1 and 2, and
\[
                    d(C_f,(Y,Y^\vee))=d(C_s,(Y',(Y')^\vee)).
\]
If the finite evaluation has rank two at a surplus-dimension fiber and is nonzero on the primitive generic class, the switch lowers the fiber dimension by two. A sequence ending at dimension one gives a nonnegative determinant valuation whenever the classes and exact self-dualities are integral at the test DVR.
\end{lemma}

\begin{proof}
We have \(H^2(C_-)=k_{\rm gen}Y^*\oplus k_{\rm gen}(Y')^*\) using upper/lower duality over the fraction field \(k_{\rm gen}\). The finite coordinate boundary pairs to zero against \(Y\); take a complement \(x\) of \(v_0\) mapping to \((Y')^*\). Then \(v_0\wedge x\) has unit plane volume by the cross pairing and the local identity. Likewise on the transverse plane interchanging the two classes. Thus the two triangles compare the class-functional tensors to the same lower tensor \((Y^*\wedge (Y')^*)^*\) times unit plane volumes. 

This comparison requires compatible duality adjunctions, as in the preceding diagram discussion. At a split-in-cohomology switch place with conditions in degrees 0,1 they may be used simultaneously: isotropic restrictions admit compatible nullhomotopies on the lower comparisons (there is no ambiguity of a morphism into scalars in degree three on those tensor products). Unchanged local conditions retain the same choices. The comparison also permits volume changes from disjoint rational local modifications when these are defined by compatible triangles.

If at a surplus-dimension fiber one arranges rank two on the old finite plane (including a nonzero evaluation of the primitive reduction), the upper localization image there is the finite plane itself. Thus the switch drops \(h=\dim H^1(C_f)\) by two, while generically the two classes give the determinant comparison just stated. 

Iteration to \(h=1\) gives nonnegative valuation for the class-functional tensor over the DVR if exact self-duality and integral classes hold there: at that fiber minimality and generic dimensions give two free lines with zero differential, paired integrally perfectly. In these switches a primitive reduction exists in a two-term DVR model before switching whenever there is a generic line: its kernel is saturated and reduction gives a nonzero fiber class.

\end{proof}

\paragraph{The local switch coordinates.}
The cross calculations will use, for example, tame evaluations at inert auxiliary primes whose residue cardinalities over a quadratic field tend to 1, with trivial local action to increasing precision. Degree one then splits by evaluation on a residue Frobenius lift and a tame generator, using the inertia/residue cochains; the two axes are isotropic for a trivial-limit scalar cup. For the inertia square cup one can retain enough roots-of-unity precision: the pure-inertia additive character at any fixed precision (Frobenius value zero) factors through cyclic inertia of much greater 2-power order since the Frobenius conjugating power is 1 to that greater precision. The cyclic scalar \(H^2\) products inflate to zero there. The mixed scalar cup is a unit with dual local bases. Conjugate transport via an inert rational Frobenius preserves axes using its square as residue lift, multiplying the inertia-generator argument by a unit; thus inserts its Tate transport on the other plane. We give the action hypotheses where applying this. 

A related pure field parity computation is for two Lagrangian conditions whose sum modulo intersection is a symmetric split plane in degree one (characteristic zero). The same upper/lower adjunction makes the upper image a self-annihilating line there, necessarily one of the two axes; switching between them thus changes \(\dim H^1\) by one in parity.

\subsection{Central Kummer lattices and local measures}

\paragraph{Kummer conditions and their duality.}
At a finite place \(v\), for an elliptic curve with coefficient \(T=T_2E\), use the local complex \(U_v=E(F_v)^\wedge_2[-1]\) when imposing the full integral ordinary Kummer condition. It maps to \(L_v\) by the Kummer map and lower truncation (integral \(H^0(F_v,T)=0\)). Its derived reductions have image all \(E[2^n](F_v)\) in degree zero and exact finite Kummer in degree one, injectively: this follows by ordinary Kummer and the Bockstein on torsion invariants, taking them to their integral Kummer classes. 

These are exact orthogonal complexes under local duality. The cup restriction is nullhomotopic over \(\mathbb Z_2\) (morphisms from the derived Kummer tensor to scalars in degree two are detected on its \(H^2\), and the Kummer pairings vanish modulo all \(2^n\)). After a nullhomotopy the local quotient duality map is an equivalence, checked by finite Kummer orthogonality and the invariants and their duals at the residue field. 

On compact cohomology the familiar sequence reads
\[
 0\longrightarrow E(F_v)^\wedge_2\longrightarrow H^1(F_v,T)
           \longrightarrow\operatorname{Hom}(E^\vee(F_v)^\wedge_2,\mathbb Z_2)
                      \longrightarrow0,
\]
by finite duality. Thus at odd finite places the compact \(H^1\) is entirely Kummer (as also at real places by the ordinary Kummer sequence, since Weil--Châtelet there has exponent two). 

The singular free coordinate at a rational dyadic place after inverting 2 uses Bloch--Kato local duality and its abelian variety comparison~\cite[Sections~3.7--3.11]{BlochKato90}: finite Kummer is the log line, and cup against it on the Tate dual is evaluation of \(\exp^*\) against the Kummer logarithm. Thus for \(\exp^*z=\alpha\omega_E\) at \(\mathbb Q_2\) the invariant against a Kummer class \(x\) is \(\pm\alpha\log_{\omega_E}x\). More generally use the differential/tangent and trace pairing. This is a rational comparison at the local prime, not a claim that \(\alpha\) is an integral quotient generator.

\begin{lemma}[The global Kummer volume]\label{lem:foundations-kummer}
For the actual all-finite Kummer complex \(C_{\rm Kum}\) over an imaginary quadratic \(K\), with enough support, derived base change to finite coefficients gives the usual Selmer in degree one by these assertions, and likewise passage to divisible coefficients. If \(\Sha(E/K)\) is finite, degree one integrally is the compact Mordell--Weil module, degree-two torsion has length \(s_K=v_2(\#\Sha(E/K))\) (the residual finite quotient in the change to \(\mathbb Q_2/\mathbb Z_2\)), and degree-three torsion has length \(\tau_g=v_2(\#E(K)_{\rm tors})\) by self-duality. The free parts in degrees one and two pair unimodularly. Thus for rank one, with \(n_P\) the index valuation on the free line of a global point \(P\),
\[
               d(C_{\rm Kum},(P,P^\vee))=2n_P+2\tau_g-s_K.          \tag{Kum}\label{eq:Kum}
\]
These valuations also use conjugate-field duality if indicated (the self-duality remains integral).
\end{lemma}

\begin{proof}
The Kummer and duality descriptions above give torsion lengths \(\tau_g,s_K,\tau_g\) in degrees one, two, and three. Their signed sum in the inverse determinant convention is \(2\tau_g-s_K\). The free point and its pairing functional each contribute its index valuation \(n_P\). This gives \eqref{eq:Kum}.
\end{proof}

\paragraph{Local measure factors.}
Write \(P_\ell(1)=L_\ell(E,1)^{-1}\). For a rational minimal differential at finite primes put \(\tau_\ell=v_2(\#E(\mathbb Q_\ell)[2^\infty])\) and \(\log_{\omega_E}E(\mathbb Q_2)=2^l\mathbb Z_2\). Haar measures give
\[
 \tau_\ell=v_2(c_\ell P_\ell(1))\quad(\ell\ne2),\qquad
 \tau_2-l=v_2(c_2 P_2(1)).
\]
Indeed the connected-point measure is the nonsingular residue size divided by \(\ell\) (the small kernel from the formal group), giving \(P_\ell(1)\); at 2 use the formal-log lattice with torsion kernel. 

At \(\mathbb R\), \(H^1(\mathbb R,T)=E(\mathbb R)^\wedge_2\) records the component factor, and a primitive Betti invariant cycle \(a\) has absolute period \(\Omega_0\) equal to the connected period, so that \(\Omega_E=\#\pi_0 E(\mathbb R)\cdot\Omega_0\).

\subsection{Bounded-series specialization and division tests}

We record four specialization principles used with these volumes. All series in these principles are bounded: their coefficients lie in a fixed complete DVR, or in its fraction field with one common denominator. This condition is stronger than belonging to the unrestricted formal series ring over the fraction field.

\begin{remark}[Specialization of determinant coordinates]\label{rem:foundations-pivots}
If a generic determinant coordinate is defined by concentration in given degrees (in particular a free degree-one cohomology determinant), at a prime where the corresponding field dimensions are the same, Gaussian pivot comparison/minimality over the local ring permits specializing the coordinate identity, even if a previous set of rational pivots vanished. In stage comparisons one may keep the pivot minors nonzero by convergence to a \emph{fixed} finite specialization, and if the resulting ranks agree, numerator, denominator and class coordinates converge by their matrix formulas. Conversely bounded torsion-length and rank tests at scalar specializations (DVR constants) keep pivots of the corresponding size nonzero in the limit.
\end{remark}

\begin{remark}[Scalar limits of measures]\label{rem:foundations-measures}
Measures can have coefficients in \(\mathcal V=W(\overline{\mathbb F}_2)\) and the constant limits can then be taken in \(\mathcal V^*=\varprojlim_n\prod_{\mathscr U}(\mathcal V/2^n)\) for the fixed ultrafilter \(\mathscr U\). This is still a complete integer DVR (nonzero elements have a finite first nonzero precision and divide uniquely by that many uniformizer factors to units). Fixed finite extensions can be embedded and used similarly with their ramification. Uniformly bounded denominators disappear after inverting 2 there. Bounded group-ring values at cyclic orders tending to infinite 2-power depth give bounded power series by reduction, \(1+u\) being a chosen generator variable. Substituting a fixed finite character (with scalar enlargement), or passing to a fixed finite group quotient, commutes with these limits. Arithmetic model matrices themselves use the smaller compact rings before further base change.
\end{remark}

\begin{lemma}[Tame division from character tests]\label{lem:foundations-series}
In a DVR-integer bounded power-series ring in \(t,\mathbf u\), a divisor \(L\) with nonzero reduction at uniformizer \(=t=0\) becomes distinguished up to a unit in one tame variable after a tame group-power coordinate change. The change may retain any fixed finite list of characteristic-zero generic ranks and nonvanishings at \(t=0\). If a bounded series \(B\) is divisible by \(L\) after inverting two at every sufficiently high two-power character in \(t\), then \(L\mid B\) integrally.
\end{lemma}

\begin{proof}
Specialize a line \(1+u_j=(1+x)^{n_j}\), one \(n_j=1\) and others successively very highly 2-divisible to separate a lexicographically leading monomial on reduction (assign the dominant lex coordinates largest \(2^{v_2(n_j)}\) in its weighted order, large relative to the lower-coordinate orders of the chosen monomial). This works simultaneously for a finite list by using the product. The line extends to a coordinate change. To retain the indicated characteristic-zero generic ranks and nonvanishings at \(t=0\), in formal log coordinates use leading homogeneous polynomials, choosing the free integers \(n_j\) avoiding finitely many proper conditions compatibly with the prescribed valuations. The same regularity test applies over a finite local coefficient order by its closed-point reduction. If a bounded \(B\) is divisible by \(L\) after inverting 2 at every sufficiently high \(2\)-power character in \(t\), then Weierstrass division in that tame variable proves \(L\mid B\) integrally. At those specializations the divisor remains regular of the same degree (evaluation always means the parameter is generator value minus one); the monic polynomial quotient has no 2-torsion, so the integral remainder vanishes there. Its bounded coefficients vanish identically by one-variable Weierstrass. One may check vanishing componentwise where the coefficient order embeds in a product of DVRs preserving the residual regularity. In general bounded nonzero series conditions in several variables can be avoided by finite-order tests by iterating the one-variable assertion (also with remaining unspecialized variables). For a nonzero bounded one-variable series the valuation at high 2-power characters is eventually the minimum coefficient valuation plus the first degree attaining it times the valuation of the parameter, by unique dominance.
\end{proof}

\begin{lemma}[The pole test]\label{lem:foundations-poles}
Let \(\mathcal O\) be a complete dyadic integer DVR. Suppose
\(f_i,g_i\in\mathcal O[[t]]\) converge coefficientwise along the fixed
ultrafilter to \(f,g\in\mathcal O[[t]]\), with \(f\ne0\). Assume, on a
filter-large set of stages, that \(f_i\ne0\) and
\[
                         g_i/f_i\in\mathcal O[[t]][1/2].
\]
No uniform bound on the denominators of these quotients is required.
Every zero of \(f\) in the open unit disk, taken in an algebraic closure
of \(\operatorname{Frac}\mathcal O\), is a zero of \(g\) with at least
the same multiplicity.
\end{lemma}

\begin{proof}
The common integral coefficient bound makes coefficientwise convergence
uniform on every smaller closed disk. Let \(\alpha\) be a zero of \(f\)
of multiplicity \(m\), and pass to a finite scalar extension containing
\(\alpha\). If \(g=0\), there is nothing to prove. Otherwise choose,
after a further finite extension if needed, a sufficiently small closed
disk about \(\alpha\) containing no
other zero of \(f\) or \(g\), and with neither series vanishing on its
boundary. After translation and rescaling, the dominant coefficient
computes the Weierstrass degree on that disk. Uniform convergence
preserves this degree for both series on a filter-large set. Thus \(f_i\)
has exactly \(m\) zeros there, while the number of zeros of \(g_i\) equals
the multiplicity of \(g\) at \(\alpha\). Stage divisibility forces the
latter number to be at least \(m\), proving the assertion.
\end{proof}

We apply this lemma after fixing the tame characters. If a denominator
factor is distinguished in \(t\), all its zeros lie in the open unit
disk. The lemma therefore makes the integral division remainder by each
required power of that factor vanish at every admissible tame-character
test. When the tests exclude only finitely many nonidentical vanishing
conditions, the bounded-series uniqueness argument above makes the
remainder vanish identically.
